\documentclass[10pt, aspectratio=169]{beamer}
\usetheme{Madrid}
\usecolortheme{default}

% Packages
\usepackage[utf8]{inputenc}
\usepackage{graphicx}
\usepackage{booktabs}
\usepackage{hyperref}
\usepackage{amsmath}
\usepackage{tcolorbox}

% Title Info
\title[MOMENT vs RmGPT]{Architectural Differences and Unsupervised Pretraining: \\ A Comparative Analysis of MOMENT and RmGPT}
\author{Robot Fault Classification Project}
\date{\today}

\begin{document}

\begin{frame}[plain]
    \titlepage
\end{frame}

\begin{frame}
    \frametitle{Objective}
    \begin{itemize}
        \item Analyze in-depth the structural differences between \textbf{MOMENT} and \textbf{RmGPT}.
        \item Detail the \textit{inductive biases} inherent in each architecture and how they affect time-series modeling.
        \item Explain the mathematical impossibility of directly comparing the \textit{Loss} values of both models.
        \item Present a \textbf{pipeline-agnostic} comparative methodology.
        \item Examine the Unsupervised Masked Patch Reconstruction experiments and the real-world dataset bottleneck (190 samples).
    \end{itemize}
\end{frame}

\begin{frame}
    \frametitle{Deep Architectural Differences: Inductive Biases}
    While both models have similar capacities ($\sim$11M parameters for MOMENT vs $\sim$12M for RmGPT), their core handling of multidimensional signals drastically diverges:
    \vspace{0.3cm}
    \begin{columns}[T]
        \begin{column}{0.48\textwidth}
            \textbf{RmGPT (Variable-Channel Bias)}
            \begin{itemize} % added detail on explicit handling
                \item Intrinsic architectural capability to separate distinct physical channels.
                \item Understands time-series features independently, preserving frequency and temporal correlations across distinct robot motors naturally.
            \end{itemize}
        \end{column}
        \hfill
        \begin{column}{0.48\textwidth}
            \textbf{MOMENT (Text-Sequence Transformer)}
            \begin{itemize} % detailed T5 flattening behavior
                \item Built strictly upon a discrete language model logic (FLAN-T5).
                \item Flattens all physical channels into continuous uniform patches, treating them as abstract ``text tokens''.
                \item Lacks native channel separation biases, meaning it must implicitly deduce structural physics from scratch.
            \end{itemize}
        \end{column}
    \end{columns}
\end{frame}

\begin{frame}
    \frametitle{The Loss Discrepancy: Why Direct Comparison is Invalid}
    In our experiments, the raw validation losses output by the trainers ($\sim 0.14$ vs $\sim 0.0001$) cannot be compared due to the \textbf{target reconstruction space}:
    \vspace{0.3cm}
    \begin{columns}[T]
        \begin{column}{0.48\textwidth}
            \textbf{RmGPT: Latent Space Loss}
            \begin{itemize}
                \item Evaluates Mean Squared Error (MSE) natively over internal mathematical token embeddings.
                \item Consequently outputs minuscule absolute error magnitudes ($<10^{-4}$).
            \end{itemize}
        \end{column}
        \hfill
        \begin{column}{0.48\textwidth}
            \textbf{MOMENT: Physical Space Loss}
            \begin{itemize}
                \item Obligatorily passes representations through its RevIN decoder before evaluation.
                \item Calculates error directly against the \textbf{normalized physical XYZ coordinates}.
            \end{itemize}
        \end{column}
    \end{columns}
\end{frame}

\begin{frame}
    \frametitle{Pipeline-Agnostic Experimental Setup}
    To eliminate external variables and ensure differences originated solely from the model architecture, we aligned the training conditions structurally (\textbf{pipeline-agnostic} framework). 
    
    \vspace{0.2cm}
    We exclusively leveraged \textbf{Unsupervised Masked Patch Reconstruction}. By hiding portions of the signal, the network is forced to logically deduce the structural physics of the robot without being biased by extreme label scarcity. Both training and validation losses converged stably in parallel.
\end{frame}

\begin{frame}
    \frametitle{Pretraining vs Fine-Tuning Convergence}
    We studied how the intensity of the pretraining phase directly dictates the fine-tuning boundary:
    
    \vspace{0.3cm}
    \textbf{1. Original Pretraining Configuration} (Batch = 32)
    \begin{itemize}
        \item Higher update frequency resulted in a deeper theoretical convergence: \textbf{Loss reached 0.14}.
        \item \textit{Fine-Tuning Impact:} Allowed the network to reach our maximal upper limit of \textbf{44\% Validation Accuracy} in real-world fault classification.
    \end{itemize}
    \vspace{0.2cm}
    \textbf{2. Adapted Configuration} (Batch = 256)
    \begin{itemize}
        \item Using slower, broader steps to mimic alternate pipelines established a plateau: \textbf{Loss reached 0.32}.
        \item \textit{Fine-Tuning Impact:} The looser representations degraded significantly, plummeting classification accuracy to \textbf{22\%}.
    \end{itemize}
\end{frame}

\begin{frame}
    \frametitle{The Real-World Data Bottleneck}
    Despite optimal unfreezing strategies and data slicing techniques, MOMENT consistently overfits reaching a hard structural ceiling:
    \vspace{0.2cm}
    \begin{tcolorbox}[colback=red!5!white,colframe=red!75!black,title=The 190 Samples Limitation vs Inference Bias]
    Because MOMENT treats data like a text transformer devoid of innate physical channel-separation biases, it mathematically requires \textbf{massive volumes of examples} to learn the structural rules of multi-motor frequency interactions. With only \textbf{190 real-world samples} distributed among 9 fault classes, the network overfits to simulation data geometries, rendering it incapable of surpassing the 44\% fine-tuning threshold.
    \end{tcolorbox}
\end{frame}

\begin{frame}
    \frametitle{Conclusions}
    \begin{enumerate}
        \item \textbf{Architectural Inductive Biases}: RmGPT's innate signal-separation handles time-series far more efficiently than MOMENT's text-token homogenization when data is scarce.
        \item \textbf{Loss Topologies}: Metrics are mathematically incomparable due to Latent vs. Physical Space evaluations. We must evaluate solely via subsequent fine-tuning behavior.
        \item \textbf{Optimization Integrity}: A deeper pretraining (Original config at $0.14$ loss) vastly outperforms looser variants ($0.32$ loss) for our specific 9-class mapping.
        \item \textbf{The Data Barrier}: Surpassing 44\% accuracy on MOMENT stringently requires exponentially expanding the empirical dataset to compensate for its text-like generalization mechanism.
    \end{enumerate}
\end{frame}

\end{document}
